\chapter{Discussion and Conclusion} \label{Chap7}

\section{Main Findings}

\subsection{Overview of Model Performance} \label{sec:model_performance_overview}

\begin{itemize}
    \item \textbf{Gap-filling:} TabICLv2 is best (combined hourly $R^2=0.515$,
    Chapter~\ref{Chap4}), leading at T2/T4 and tying RF at T9. The two neural models nevertheless 
    underperformed compared to the leading tree-based and foundation-model approaches:
    bidirectional LSTM ($R^2=0.224$) and SAITS ($R^2=0.385$).
    \item \textbf{Forecasting:} TabPFN is best (MASE$=0.6908$, Chapter~\ref{Chap5}), while TabICLv2
    ranks second (MASE$=0.7191$). The weakest were again the locally trained deep-learning models:
    DLinear (MASE$=1.1872$, the only model worse than the seasonal-mean baseline) and LSTM (MASE$=0.9521$).
    \item \textbf{Tabular foundation models excel.} TabPFN and TabICLv2 are pretrained on
    millions of synthetic tabular datasets and predict via in-context learning, fitting no
    dataset-specific parameters \cite{hollmann2023tabpfn,hollmann2025tabpfn,qu2025tabicl,qu2026tabiclv2}.
    This prior-fitted-network approach is specifically effective on small, incomplete datasets
    \cite{hollmann2025tabpfn} and may provide an advantage on NWFP's severely incomplete record.
    \item \textbf{Deep learning does not perform best.} Locally trained neural architectures, 
    including LSTM, Bi-LSTM, DLinear and TFT train from scratch on NWFP's own incomplete 
    record, giving a small effective training-sample size. Off-the-shelf recurrent architectures 
    without dedicated missingness-handling tend only to match classical baselines (RF, SVM) on incomplete
    time series rather than exceed them \cite{che2018grud}. This is consistent with deep learning's broader
    failure to reliably beat tree-based methods on tabular data \cite{grinsztajn2022}.
    \item \textbf{Uncertainty quantification also favours the tabular foundation models.} In
    gap-filling, TabICLv2 matches QRF's calibrated coverage (PICP$=0.899$) with sharper intervals
    (MPIW $176.2$ vs.\ $183.5$, Chapter~\ref{Chap4}). In forecasting, TabPFN's raw intervals
    already sit close to nominal at T4/T9 (PICP$=0.929$/$0.928$) without extensive calibration --
    though not at Tower 2, where too few observations remain to calibrate at all
    (Chapter~\ref{Chap5}).
\end{itemize}

\subsection{Difficulty in Methane Flux Gap-Filling and Forecasting}

\begin{itemize}
    \item \textbf{Gap-filling fit remains modest despite the improved pipeline.} Even with
    enriched features and a tabular foundation model, the combined hourly $R^2$ reaches only
    0.515 (Chapter~\ref{Chap4}), a large improvement over MDS ($R^2=0.055$), but still a modest 
    absolute fit.
    \item \textbf{The livestock-dominant catchments are hardest to reconstruct.} T4 and T9, where
    livestock density dominates the FCH\textsubscript{4} signal, score \emph{below} the combined
    average ($R^2=0.429$/$0.430$), while T2, driven by a land-use transition, scores highest ($R^2=0.686$). 
    The catchments most central to this dissertation's finding are the ones it reconstructs least well.
    \item \textbf{Forecasting error is dominated by spikes.} Section~\ref{fc:spike} shows the
    champion's strong pooled MASE (0.6908) is earned almost entirely on ordinary days (non-spike
    MASE$=0.524$); on spike days it rises to 3.285, over six times worse and a clear loss against
    baseline. No feature in this project reliably predicts spike occurrence.
\end{itemize}

  \subsection{Scenario Sensitivity}

  Livestock density produced the largest modelled response among the tested
  scenarios, with historical-maximum stocking levels increasing predicted annual
  FCH\textsubscript{4} by 100--184\%. Extending grazing increased predictions by
  approximately 20\% at T4 and T9, whereas fertiliser interventions produced
  smaller and less consistent responses. These sensitivities were similar under
  SSP2-4.5 and SSP5-8.5, although they represent conditional model behaviour
  rather than validated causal effects.

\section{Limitations} \label{sec:limitations}

\begin{itemize}
    \item \textbf{Tower 2 is the weakest-conditioned data source at every pipeline stage.} Its
    QC-valid coverage is lowest of the three towers (92.0\% unavailable vs.\ 68.3\%/81.7\% at
    T4/T9, Chapter~\ref{Chap3}); it cannot support conformal calibration in forecasting
    (Chapter~\ref{Chap5}) due to data limitations.
    \item \textbf{Forecasting and scenario projection rely substantially on gap-filled variables.} Both
    build on Chapter~\ref{Chap4}'s reconstructed FCH\textsubscript{4} series, most of which is
    model-estimated, not observed (68.3--92.0\% of hours). Downstream results therefore partly
    reflect the gap-filling model's own behaviour.
    \item \textbf{UQ and scenario projection were not extended to the full model roster.}
    Gap-filling UQ covers only its top two models (QRF, TabICLv2, Chapter~\ref{Chap4});
    forecasting UQ covers only its champion (TabPFN, Chapter~\ref{Chap5}) -- extending further was
    not attempted given the compute cost. 
    \item \textbf{CMIP6 climate simulations do not match NWFP's own historical climate.} Raw air
    temperature runs 1.81--2.71$^\circ$C too cool and precipitation $\sim$4.2--4.4$\times$ too wet
    before correction (Chapter~\ref{Chap6}); the applied correction reduces but does not eliminate
    this as a residual uncertainty source.
    \item \textbf{Scenario projection has no ground truth to validate against.} No real
    CH\textsubscript{4} measurement exists for conditions that have never occurred at NWFP, and
    this dissertation makes no comparison against any process-based (mechanistic) methane model; 
    its results describe conditional model behaviour, not a validated forecast.
\end{itemize}

\section{Future Work}

\begin{itemize}
    \item \textbf{Rerun the experiment once extended Tower 2 data is released.} Tower 2's usable
    record is expected to grow soon: 2025 data has already been collected and is pending release.
    Rerunning once available would test whether Tower 2's results reflect a resolvable
    data-scarcity limitation (Section~\ref{sec:limitations}) rather than a persistent structural
    difference.
    \item \textbf{Livestock proximity to each tower's footprint as an additional gap-filling
    feature.} This dissertation's livestock-density feature is catchment-level, but an EC
    footprint shifts with wind direction, speed, and stability (Chapter~\ref{Chap2}), so
    catchment-level presence does not guarantee animals were within it. Proximity-weighted,
    footprint-aware livestock location via GPS or remote sensing could sharpen the single
    strongest predictor identified throughout this dissertation.
    \item \textbf{Benchmark tabular models against other eddy-covariance datasets.} This
    dissertation's finding that pretrained tabular foundation models outperform locally-trained
    alternatives under severe missingness (Section~\ref{sec:model_performance_overview}) was
    established at a single site. Testing the same roster against benchmarks such as FLUXNET-CH4
    \cite{irvin2021} would test whether this generalises. However, FLUXNET-CH4 predominantly
    covers wetlands and rice paddies rather than managed grasslands, so this would specifically test
    the pretrained-model and environmental-driver findings.
\end{itemize}

\section{Concluding Remarks}
Overall, this dissertation developed and evaluated an integrated framework for
methane-flux gap-filling, conditional forecasting and exploratory scenario
analysis at the North Wyke Farm Platform. Pretrained foundation models for
tabular data provided the strongest overall predictive performance, while
traditional tree-based methods generally outperformed locally trained neural
networks. Livestock-associated variables also formed the most consistent
influence across the forecasting and scenario experiments. These findings
remain constrained by substantial target missingness, poorly predicted
emission peaks and the absence of future ground truth, and the scenario
responses should not be interpreted as causal effects. Nevertheless, the
resulting benchmark provides a reproducible foundation for further methane
modelling at North Wyke and other managed grassland sites.
